% ==============================================================================
% 04-nonfunctional.tex — Раздел 2.3. Нефункциональные требования (сокращённо)
% ==============================================================================

\section{Нефункциональные требования}
\label{sec:nonfunctional}

НФТ структурированы по модели качества ISO/IEC~25010~\cite{src:9, src:10} (таблица~\ref{tab:nfr}).

\begin{table}[H]
  \caption{Нефункциональные требования (ISO/IEC 25010)}
  \label{tab:nfr}
  \small
  \begin{tabularx}{\textwidth}{|p{3.8cm}|X|}
    \hline
    \textbf{Характеристика} & \textbf{Требование} \\
    \hline
    Функциональная пригодность & Покрытие сценариев: каналы, консистентность, обратная связь \\
    \hline
    Производительность & Приемлемое время операций при росте каналов и сущностей \\
    \hline
    Совместимость & Интеграции через программный интерфейс или браузерную автоматизацию, единый интерфейс агента, расширяемость платформ \\
    \hline
    Удобство & Поддержка ролей, минимизация ошибок при массовых обновлениях \\
    \hline
    Надёжность & Целостность данных и история при частичных сбоях интеграций \\
    \hline
    Безопасность & Шифрование токенов, аудит, разграничение доступа \\
    \hline
    Сопровождаемость & Добавление агентов и изменение оркестрации без переработки ядра \\
    \hline
    Переносимость & Развёртывание в~контейнерах \\
    \hline
  \end{tabularx}
\end{table}

\textbf{Целевые показатели:} 95-й перцентиль времени отклика интерфейса < 500~мс, ответ языковой модели < 10~с, синхронизация по платформе < 30~с, доступность > 99{,}5\,\%, время восстановления < 1~ч. Сервер: 2~ядра, 4~ГБ оперативной памяти, 20~ГБ дискового пространства; дополнительно не менее 1~ГБ на агент с~браузерной автоматизацией. Круглосуточный режим работы, браузеры~--- последние две стабильные версии. Документация: описание интерфейса по стандарту OpenAPI, инструкция по развёртыванию, руководство пользователя; техническое задание по ГОСТ~34.602-2020~--- в~Приложении А к~пояснительной записке ВКР.
